\documentclass{article}

\PassOptionsToPackage{numbers,sort&compress}{natbib}
\usepackage[final]{neurips_2019}

\usepackage[utf8]{inputenc}
\usepackage[T1]{fontenc}
\usepackage{hyperref}
\usepackage{url}
\usepackage{booktabs}
\usepackage{multirow}
\usepackage{amsfonts}
\usepackage{nicefrac}
\usepackage{microtype}
\usepackage{graphicx}
\usepackage{xcolor}
\usepackage{lipsum}
\usepackage{fvextra}
\usepackage{tikz}
\usetikzlibrary{arrows.meta, positioning, fit, backgrounds, decorations.pathreplacing, calc}
\usepackage{fontawesome5}
\usepackage{amsmath}
\usepackage{cleveref}
\graphicspath{{report/}}
\newcommand{\yx}[1]{\textco/home/ramneek/Documents/QuantBeyoundEfficiency/report/acl_natbib.bstlor{blue}{{[YX: #1]}}}

\title{
  Responsible AI: LM Safety Harness \\
  \vspace{1em}
  \small{\normalfont Stanford CS224N \{Custom, Default\} Project}  % Select one and delete the other
}

\author{
  Hans \\
  Department of Computer Science \\
  Stanford University \\
  \texttt{hans@stanford.edu} \\
  \And
  Ramneek \\
  Department of Computer Science \\
  Stanford University \\
  \texttt{ramneek@stanford.edu} \\
}


\begin{document}

\maketitle

\begin{abstract}
	
\end{abstract}

\section{Introduction}
Large language models are becoming increasingly capable of performing safety-critical tasks, such as those in the fields of cybersecurity, chemistry and biology.
However, even in state-of-the-art models such as Claude Fable 5, jailbreaking remains a significant attack vector for language models. 
This refers to the model's safeguards being bypassed by attacks carried out within the context window by the user. 
Given their increasing capabilities, this remains a dangerous weakness of LMs. 
In this work, we treat the LM as a safety-critical system and aim to mitigate the risks of LM misuse by creating a safety harness around the LLM.
A safety harness is defined as a software scaffold around an LLM. Unlike agent harnesses, where the scaffolds are RAGs or connections to tool calls, our method connects LLMs to various detectors. We use BERT classifiers to detect input and output text, while simultaneously auditing the models' internals using linear probes.
We evaluated the models using the following metrics: F1 score, accuracy and latency, comparing different combinations of detectors. 
The final decision on whether an input or generation is flagged as a jailbreak is made by combining the majority vote from all the classifiers. 
\par 

In this work we ask the following research questions: 
\begin{itemize}
    \item \textbf{RQ1:} How robust are safety classifiers for LLMs when combined? 
    \item \textbf{RQ2:} Can we successfully implement internal methods during runtime without inhibiting the model's capabilities? 
    \item \textbf{RQ3:} What are the limitations of current LLM safeguards?  
\end{itemize}

These are our contributions: 
\begin{itemize}
	\item In this work we created a flexible pipeline which allows interchangable and flexible usage of different safety mechanisms.
    \item Our architecture audits the input, internals and the outputs of the model simultaneously  
	\item  We have implemented an architecture which combines a finetuned input bert model with multiple linear probes and a finetuned output bert model, 
to detect jailbreaking attacks above baseline. 

\end{itemize}

\section{Background}


\section{Methodology}

\subsection{Software-Architecture}
\begin{figure}[!ht]
\centering
\includegraphics[width=1\textwidth]{figures/architecture_basic}
\caption{title}
\end{figure}
See the LLM as a safety-critical system and audit the input, internals, and output in real time.
Potential solutions include:

- Input: Safety Classifier to detect malicious prompts or jailbreaking attempts
- Prompt rewriting
- Internals: Sparse autoencoders, linear probes, and steering vectors. 
Output: Safety Classifier 

\subsection{Bert Model}

\subsection{Linear Probe}
\subsubsection{Input Linear Probe}
A linear probe is a linear decision boundary trained using the internal activations of a large language model (LM).
The first linear probe was trained using the final token activation from the residual layer of the LM. 
According to the linear representation hypothesis, features are approximately linear directions in the activation space. 
Therefore, if the hyperplane is approximately orthogonal to the jailbreaking feature, a linear probe can sufficiently classify the activations as either negative or positive.  
We trained the linear probe using a dataset with five epochs. 
\subsubsection{Predictive Linear Probe}
This probe predicts whether the model input would result in a jailbreak. 
Therefore, it calculates the probability of a jailbreak based on the assumption 
that the input is already a jailbreak (conditional probability). 
The model was trained using a dataset consisting only of jailbreak attempts, 
which were then passed through the model and evaluated using an LLM judge.
 The linear probe was trained in the final residual layer of the language model with the judged labels. 
\subsection{Evaluation}
Finally, evaluate combinations and different methods using an evaluation or red teaming method. 
Evaluation Methods: 
Standard metrics: F1 Score, Accuracy, latency, Confusion Matrix. 


\section{Experiments}

\bibliographystyle{unsrtnat}
\bibliography{references}

\end{document}